\chapter{Breast Reduction}

\section*{Definition and Overview}
Breast reduction, also called reduction mammoplasty, is a surgical procedure that removes excess breast tissue, fat, and skin to make the breasts smaller, lighter, and firmer, and to reposition the nipple and areola to a higher, more natural position. The operation is performed primarily to relieve the physical symptoms caused by large, heavy breasts -- a condition known as macromastia -- and, in most cases, to improve the appearance and comfort of the breasts. For many women, surgery is the only intervention that provides lasting relief from the pain, skin problems, and functional limitations of very large breasts. The procedure is usually performed under general anaesthesia and typically uses a vertical or an anchor (inverted-T) incision pattern, although the exact technique depends on the size, shape, and ptosis (sagging) of the breasts, the amount of tissue to be removed, and the surgeon's approach.

\section*{Epidemiology}
Breast reduction is one of the more commonly performed plastic surgery procedures in women, although it is far less common in men. Women with macromastia come from all age groups, but surgery is most often requested in early adulthood and middle age. Approximately 1--5\% of women have breasts large enough to cause significant physical symptoms, and the condition tends to run in families. Large breast size is commonly associated with being overweight, and weight gain enlarges the breasts because they contain a high proportion of fatty tissue. Men may also develop excess breast glandular tissue (gynaecomastia), for which a smaller reduction-type procedure may be offered. Access to publicly funded surgery varies between health systems and usually depends on the severity of symptoms and local funding guidelines.

\section*{Aetiology and Pathophysiology}
\begin{itemize}
    \item \textbf{Genetics and hormones:} breast size is strongly influenced by inherited factors and by the female sex hormones oestrogen and progesterone during puberty, pregnancy, and weight change
    \item \textbf{Weight gain:} the breasts contain a large proportion of fatty tissue, so increases in body weight commonly enlarge them
    \item \textbf{Gynaecomastia in men:} an excess of glandular breast tissue, often related to hormonal imbalance, certain medicines, or liver and thyroid conditions
    \item \textbf{Mechanical effects:} heavy breasts pull on the skin and soft tissues of the chest, shoulders, and back, straining the supporting ligaments and altering the posture of the neck and upper spine
    \item \textbf{Physical consequences:} the weight of large breasts causes deep bra-strap grooves in the shoulders, skin irritation and intertrigo beneath the breasts, and chronic pain in the neck, shoulders, and lower back
\end{itemize}

\section*{Clinical Features}
The main features of macromastia relate to the weight and bulk of the breasts rather than to any underlying breast disease:

\begin{itemize}
    \item Neck, shoulder, and upper or lower back pain that typically improves with a supportive bra
    \item Deep grooves or indentations in the shoulders caused by bra straps
    \item Skin irritation, chafing, sweat rash, or fungal infection beneath the breasts
    \item Numbness or tingling in the hands or fingers, sometimes related to posture and muscle strain
    \item Headaches related to chronic neck and shoulder muscle tension
    \item Difficulty finding clothes and bras that fit, and problems with exercise and physical activity
    \item Discomfort or self-consciousness that affects body image and quality of life
\end{itemize}

\section*{Differential Diagnosis}
\begin{itemize}
    \item \textbf{Physiological breast enlargement:} large breasts in a woman of otherwise healthy weight
    \item \textbf{Weight-related enlargement:} breast growth largely attributable to generalised obesity
    \item \textbf{Breast lump or tumour:} a mass causing localised enlargement or asymmetry, including breast cancer
    \item \textbf{Mastitis or infection:} a swollen, tender, or painful breast, especially during breastfeeding
    \item \textbf{Hormonal enlargement:} breast swelling related to pregnancy, the menstrual cycle, or hormone therapy
    \item \textbf{Gynaecomastia in men:} other causes of male breast enlargement, including medication and endocrine disorders
\end{itemize}

\section*{When to Seek Medical Care}
See a doctor if large breasts are causing persistent pain, skin problems, or significant difficulty with daily activities such as exercise, work, or sleep, and if these symptoms are not relieved by a well-fitted, supportive bra. As with any breast symptom, it is also important to exclude underlying breast disease before attributing changes to simple size.

\textbf{Call 000 (or local emergency services) for:}
\begin{itemize}
    \item Any new breast lump, especially one that is hard, fixed, or growing
    \item Sudden swelling, redness, heat, or fever suggesting a breast infection
    \item Bloody or unusual nipple discharge
    \item A change in the skin of the breast, such as dimpling, puckering, or a newly inverted nipple
\end{itemize}

\section*{Diagnosis and Investigations}
\begin{itemize}
    \item \textbf{History:} a detailed account of breast symptoms, weight changes, pregnancies, breastfeeding history, and current medicines
    \item \textbf{Examination:} assessment of breast size, symmetry, skin condition, and the position of the nipples, together with general weight and posture
    \item \textbf{Breast imaging:} mammography or ultrasound may be performed in older women, or where there are lumps or concerning features, to exclude breast cancer before surgery
    \item \textbf{Photographs:} taken for surgical planning, documentation, and measurement of the result
    \item \textbf{Medical assessment:} a general health review, including blood pressure and any heart or lung conditions, to ensure fitness for an operation under anaesthesia
\end{itemize}

\section*{Management}
\begin{itemize}
    \item \textbf{Non-surgical measures:} weight loss, a well-fitted supportive bra, physiotherapy for neck and back strain, and treatment of skin irritation beneath the breasts; these may help but usually do not fully relieve the symptoms of true macromastia
    \item \textbf{Breast reduction surgery:} removal of excess breast tissue, fat, and skin with repositioning of the nipple and areola; techniques include the inferior pedicle, vertical scar, and free nipple graft approaches
    \item \textbf{Liposuction-assisted reduction:} removal of fatty tissue with a suction cannula, sometimes combined with standard reduction techniques, and used alone in some cases of predominantly fatty enlargement
    \item \textbf{Post-operative care:} dressings, wound care, and a supportive bra worn for several weeks, with a gradual return to normal activity over four to six weeks
    \item \textbf{Surgeon counselling:} a detailed discussion of the expected size, the position of scars, and the potential complications before deciding to proceed
\end{itemize}

\section*{Complications}
\begin{itemize}
    \item Bleeding, bruising, or a collection of fluid or blood under the skin (haematoma or seroma) requiring drainage
    \item Infection or delayed wound healing, sometimes resulting in widened or unsightly scars
    \item Loss of feeling in the nipple or breast skin, which may be temporary or permanent
    \item Reduced ability to breastfeed in the future, because the milk ducts may be divided
    \item Changes in nipple position, and rare loss of the nipple or areola if its blood supply fails
    \item Asymmetry, or the need for further surgery to refine shape or size
    \item Scarring that takes up to a year or more to mature and fade
    \item The general risks of anaesthesia and surgery, including venous thromboembolism
\end{itemize}

\section*{Prognosis and Outcomes}
Breast reduction is one of the most satisfying procedures in plastic surgery in terms of symptom relief. The great majority of women report marked improvement in neck, back, and shoulder pain, better skin health beneath the breasts, and an improved ability to exercise and take part in physical activity, together with substantial gains in body image and confidence. Most surgical scars fade considerably over six to twelve months, although they do not disappear completely. The results are generally long-lasting, although weight gain, pregnancy, breastfeeding, and the natural ageing process can all alter breast size and shape over time. As with all surgery, careful follow-up and attention to wound care optimise the final result.

\section*{Prevention and Self-Care}
\begin{itemize}
    \item Maintain a healthy weight, since weight gain tends to enlarge the breasts
    \item Wear a well-fitted, supportive bra, especially during exercise
    \item Keep the skin beneath the breasts clean and dry to prevent rashes and fungal infection
    \item Use moisturiser or a barrier cream on skin creases where irritation develops
    \item Practise good posture and consider physiotherapy or stretching for neck and back strain
    \item If considering surgery, choose an experienced surgeon and discuss expectations, scars, and risks fully
    \item Attend all follow-up appointments and follow wound-care advice carefully after surgery
\end{itemize}

\section*{Sources and Further Reading}
The following organisations provide reliable, evidence-based information for patients, students, and clinicians.

\begin{itemize}
    \item NHS. \textit{Breast reduction: what the operation involves and recovery.} https://www.nhs.uk/conditions/
    \item Mayo Clinic. \textit{Breast reduction surgery: overview and patient education.} https://www.mayoclinic.org/diseases-conditions/
    \item MSD Manual. \textit{Breast reduction (reduction mammoplasty): indications and techniques.} https://www.msdmanuals.com/
    \item MedlinePlus. \textit{Breast reduction: consumer health information.} https://medlineplus.gov/
    \item National Institute for Health and Care Excellence. \textit{Clinical guidance on breast reduction surgery.} https://www.nice.org.uk/
    \item World Health Organization. \textit{Information on surgical services and safe surgical care.} https://www.who.int/
\end{itemize}
